\subsection{\jazzmine}

\jazzmine is the agent responsible for fuzzing Java targets using a modified version of the Jazzer tool~\cite{jazzer}.

As the first step, \jazzmine uses CodeQL to extract all of the interesting strings from the target, in order to build a dictionary, which will be used during the fuzzing process.

Then, \jazzmine retrieves from the \builder a version of the target with all the sanitizers provided natively by Jazzer and, in addition, extended custom-built sanitizers.
These sanitizers are ``loosened'' version of the original sanitizers that allow for less strict match between the values in the inputs and the values used in monitored instructions of the application.
Which sanitizers are used in a specific execution can be selected through an environment variable (choice of the original set, the extended one, or both).

% \TBD: Explain why the sanitizers have been loosened with an example.
Unlike sanitizers developed for compiled languages, Java sanitizers are highly bug class-specific and employ a fundamentally different approach. These sanitizers operate by hooking sink functions and performing validation checks on data as it reaches these critical points as function arguments. However, their implementation relies on hardcoded strings known as sentinels for data control validation. A prime example is the working of OS command injection sanitizer, which hooks the \textbf{start} function of the \textbf{java.lang.ProcessImpl} class in the JDK and subsequently verifies whether this function can be invoked with \textbf{"jazze"} as the command argument. While this approach proves effective in scenarios with straightforward Input-to-State Correspondence, it becomes significantly more challenging when compression, encoding, or other obfuscation techniques are introduced. Under these conditions, the fuzzer struggles to accurately determine and generate the precise input required to trigger the sanitizer checks, since it relies solely on fuzzer mutation and execution to build the Input-to-State Correspondence. 

To address this limitation, we introduce an LLM agent that figures out the Input-to-State Correspondence and generates the required input to reach the state with the correct sentinel.  We modify the sanitizer to trigger when the hooked function is called with fuzzer artifacts, which are strings or characters that have a high probability of being generated by the fuzzer's mutator, thereby expanding or "loosening" the conditions under which the sanitizer crashes.

Then, \jazzmine calls \grammarguy and determines if a grammar has been produced for the relevant harness.

%\TBD: What is actually done with this grammar?

This grammar generates test inputs that are validated against the original sanitizer; this is the step which is used to filter out all the false positive crashes that might have been introduced due to the loosening of the sanitizer. We use the original sanitizer's crash as the ground truth for the LLM agent, if it is able to crash, then we mark the seed as the crashing input and send it to POV guy. If the agent fails to produce the crashing input within the allocated budget or iterations, we preserve the grammar and send it to \revolver for fuzzing.

There are three possible campaign modes:

\begin{enumerate}
\item original mode;
\item with extended sanitizers;
\item with grammar-based mutators (if a grammar was generated for the harness).
\end{enumerate}

The seeds for each harness are synchronized across campaigns. 
In addition, \jazzmine minimizes the seed corpus for all the fuzzers running on a shared Kubernetes node.

\jazzmine also sync seeds across nodes, and receives seeds from \grammarguy, \quickseed, and \discoveryguy.

%We also have the "secret" seeds from OSS-Fuzz.

When a crash is found, the harness and seed is sent to \povguy.

%When a crash is based on an extended sanitizer, it is sent to \grammaragentlowsanreproducer, who is responsible for making the same crashing input trigger the corresponding standard sanitizer.













